\section{Related Work}

\subsection{Scene Graph Generation}
Scene graph generation belongs to visual relationship detection and uses a graph to represent visual objects and the relationships between them~\cite{scene-graph}. In early works, detection and prediction are performed only based on specific objects or object pairs~\cite{Lu16}, which can't make full use of the semantic information in the input image. To take the contextual information into consideration, later works utilize the whole image by designing various network architectures, such as biLSTM~\cite{Zellers_2018_CVPR}, TreeLSTM~\citep{VCTree_tang} and GNN~\cite{Yang_2018_ECCV,li2021relationship}.
%which can pass information between different objects and object pairs.
%However, these methods are all limited by the insufficient label information and imbalanced data distribution of the training dataset.
Some other works try to utilize external information, such as linguistic knowledge and knowledge graph, to further improve scene graph generation~\cite{Lu16, chen2019knowledge, Yu2017,gu2019scene}.

However, due to the annotator preference and image distribution, this task suffers from a severe data imbalance issue. Dozens of works thus try to address this issue, and 
%regard this task as unbiased scene graph generation, which is a balanced form of SGG task.
the main focus is on reducing the influence of long-tail distribution and building a balanced model for relationship prediction. Tang et al. also try to use causal analysis \cite{Tang2020_unbiased} to reduce the influence of training data distribution on the final model. Some other works \cite{chen2019soft,zhan2020multi,chiou2021recovering} address this issue in a positive-unlabeled learning manner, and typical imbalance learning methods, such as re-sampling and cost-sensitive learning, are also introduced for scene graph generation~\cite{li2021bipartite, yan2020pcpl}. Unlike these approaches, we adopt the resistance training strategy using prior bias (RTPB), which utilizes a resistance bias item for the relationship classifier during training to optimize the loss value and the classification margin of each type of relationship.

\subsection{Imbalanced Learning}

Imbalanced learning methods can be roughly divided into two categories: resampling and cost-sensitive learning.

\subsubsection{Resampling} Resampling methods change the class ratio to make the dataset a balanced one. These methods can be divided into two categories: oversampling and undersampling. Oversampling increases the count of the less frequent classes~\cite{chawla2002smote},
%Oversampling is effective in some cases but
and hence may lead to overfitting for minority classes. Undersampling reduces the sample of major categories, and thus is not feasible when data imbalance is severe since it discards a portion of valuable data~\cite{liu2008exploratory}.

% A significant disadvantage of this method is that it discards a portion of available data.
% randomly from majority classes

\subsubsection{Cost-sensitive Learning}
Cost-sensitive learning assigns different weights to samples of different categories \cite{elkan2001foundations}.
% Some methods change the weight of different samples or different classes during training.
Re-weighting is a widely used cost-sensitive strategy by using prior knowledge to balance the weights across categories. Early works use the reciprocal of their frequency or a smoothed version of the inverse square root of class frequency to adjust the training loss of different classes.
However, this tends to increase the difficulty of model optimization under extreme data imbalanced settings and large-scale scenarios. \citeauthor{cui2019class}
%argue that when the number of samples increases, the additional benefit of a newly added data point will diminish but the weight is not infinity. Therefore, they
propose class balanced (CB) weight, which is defined as the inverse effective number of samples\cite{cui2019class}. Some other approaches dynamically adjust the weight for each sample based on the model's performance~\cite{Lin_2017_ICCV, Li_Liu_Wang_2019, cao2019learning}.
%, e.g., Focal loss \cite{Lin_2017_ICCV} decreases the weights of samples that are better classified. Li et al.~\cite{Li_Liu_Wang_2019} down-weight the huge amount of easy examples and relatively down-weight the outliers. LDAM attempt to assign different margins of classification boundary for different classes \cite{cao2019learning}.

Different from these approaches, which directly deal with the loss function, we add a prior bias on the classification logits for each class based on the distribution of training dataset. In this way, we provide a new approach to import additional category-specific information and address the imbalance problem through adjusting the classification boundaries.
A similar idea is proposed in~\cite{menon2020long} for long-tailed recognition, and we differ from them significantly in the tasks and formulations.
%\par In addition to these two series of methods, recently, some other work also try to use other learning strategies such as transfer learning, metric learning, and meta-learning to solve this problem
% \TODO[relationship to us]